\section{经济与社会：把帝国写进田亩、道路与家庭}

汉代经济社会史不能只从“轻徭薄赋”“繁荣”或“土地兼并”一语开始。田租、算赋、徭役和户籍把家庭写进国家账目，却要通过乡里、县吏、地方首领和实际收获才会成为征收。逃亡、隐匿、灾荒和战争都使纸面整齐接受检验。元始二年户口与后汉帝纪中的数字，是行政看见自身的一种尺度，不是可直接换算人口、财富或家庭负担的透明统计。

钱币与市场也有同样的限度。铸钱开放和收束、五铢钱、盐铁、均输和平准，把铜料、技术、商人和官吏带进国家财政。出土钱币和钱范让这种制度有可触的物质面貌，却不能由某一窖藏推断全国货币量或市场活跃度。国家有时直接经营，有时借市场和地方中介取得资源；两者都不能离开运输与执行。

战争、筑城和边郡经营改变人群的居处。西汉向河西、岭南等地设置郡县，东汉的度田、交趾战争、羌地冲突和汉末战乱，都牵动迁徙、安置、编户与地方保护。正史多保存将领、诏令和“平定”，难以连续记录逃户、佃作、妇女、儿童和非精英家庭的选择；不应因此将他们从帝国的成本中删去，也不能在空白处虚构一份未存的社会账簿。

家族和豪强不是一个可用于一切地区的反国家标签。它们可提供粮食、保护、荐举和地方信息，也可能隐匿人户、积累部曲或把公共职位变成私人的资源。中枢依赖它们，正因中枢无法亲自抵达每一处；汉末它们同州郡和军队的结合，更使同一套地方联系能够服务不同中心。评价“国家能力”时，必须同时问命令抵达了哪里、谁承担成本、谁保留了拒绝或讨价还价的空间。

\begin{debate}[数字与遗存能说多远]
《汉书·食货志》《地理志》、后汉帝纪和简牍保留了极珍贵的数量、制度与出入文书；它们分别服务于行政申报、史书编纂或特定地点的事务。数字的口径、漏报与覆盖范围，钱币和墓葬的样本偏差，都不允许把局部材料变为全国平均数。本节据此辨认国家怎样尝试取得资源，而不制造无法核验的总产量、税率或人口曲线。
\end{debate}
